\documentclass{situationist-preprint}

\title{The Arrow of Time and the Past Hypothesis}
\msid{TSM-05-a}
\msfamily{TSM-05}
\msdiscipline{Time, Space and Modality}
\mstype{Principal manuscript}
\mskeywords{direction of time; entropy; second law; past hypothesis; statistical mechanics; best system; explanation; Boltzmann}

\begin{document}
\maketitle

\begin{abstract}
The fundamental dynamical laws of physics are, to a good approximation, invariant under time reversal, yet the world is full of processes that run only one way: heat flows from hot to cold, records exist of the past and not of the future, and we can influence what comes later but not what came earlier. Why does time have a direction? We argue that it has one because the universe began in a state of very low entropy. Given time-symmetric dynamics and a uniform probability measure over the microstates compatible with each macrostate, this past hypothesis is necessary and sufficient, in a sense we make precise, for the thermodynamic asymmetry, and the asymmetries of records, knowledge and influence derive from the thermodynamic one. We then argue that the past hypothesis requires no further explanation. It is an axiom of the best system, where explanation by subsumption ends; the charge that it is improbable misapplies a measure whose role in the system is conditional on it; and no rival of comparable standing makes it more probable.
\end{abstract}

\section{The problem}\label{s:problem}

Ice cubes melt in warm drinks and do not form spontaneously in them; perfume spreads through a room and does not gather back into the bottle; a dropped glass shatters and its fragments do not leap together. Each of these processes, filmed and played backwards, shows something that never happens. Yet the dynamical laws that govern the constituents of ice, perfume and glass, whether classical or quantum, are invariant under time reversal: for every solution of the equations of motion there is another, obtained by reversing the order of states and the velocities, that is equally a solution. The time-reversed film depicts a lawful history. The asymmetry is not in the laws.

The thermodynamic asymmetry is not the only one. We have records of the past, such as footprints, photographs and memories, and no comparable records of the future. We can affect what happens later by what we do now, and we cannot affect what happened earlier. Waves radiate outward from sources and are not seen converging on sinks. Our knowledge of the past is, in some respects, far more extensive than our knowledge of the future, which EPI-25-a argues is confined to what follows by law from present conditions. These asymmetries are pervasive and robust, and together they are what we mean when we say that time has a direction: there is a difference between the past-to-future orientation of the temporal order and the future-to-past one, and the difference shows in what happens along it.

The question of this manuscript is why time has a direction in this sense. We answer that it does because the universe began in a state of very low entropy, and that this fact, the \emph{past hypothesis}, requires no further explanation.

The question presupposes nothing about temporal passage. TSM-01 argues that time is constituted by the B-series, TSM-02 that time does not flow, and TSM-15 that the present has no metaphysical privilege. The direction of time is not a movement of the present along the B-series but a structural fact about it: the physical processes distributed along it are oriented, almost all of them the same way. The question is why the two orientations of the temporal order differ in what happens along them, given that the laws do not distinguish them.

\subsection{A history of the debate}\label{ss:history}

Boltzmann's H-theorem of 1872 purported to derive entropy increase in an isolated gas from the dynamics of colliding molecules. \citet{Loschmidt1876} objected that no such theorem could follow from time-symmetric dynamics: every entropy-increasing trajectory has an entropy-decreasing time reverse. The objection succeeded; the theorem rests on the hypothesis of molecular chaos, which is itself time-asymmetric. Boltzmann then read the second law statistically \citep{Boltzmann1896}: of the microstates compatible with a non-equilibrium macrostate, the overwhelming majority evolve toward higher entropy. But the counting is symmetric. The same majority came from higher entropy, so that, on the statistics alone, a half-melted ice cube is most probably a fluctuation out of lukewarm water, not the remnant of a larger cube \citep{Albert2000}.

Boltzmann considered two ways out. The fluctuation hypothesis, which he attributed to his assistant Schuetz, holds that the universe is in equilibrium and that our region is a rare fluctuation from it \citep{Boltzmann1895}. \citet{Eddington1928} and \citet{Feynman1965} objected that a fluctuation just large enough to produce what we now see is vastly more probable than one producing a whole ordered universe with a history; the modern form of the objection is the problem of \enquote{Boltzmann brains} \citep{DysonKlebanSusskind2002,AlbrechtSorbo2004}. A hypothesis on which our records are almost certainly illusions undermines the evidence for itself. The other way out is to assume that the universe began in a state of low entropy: the past hypothesis, so named by \citet{Albert2000}. \citet{Reichenbach1956} tried to derive the arrow from the parallel behavior of quasi-isolated \emph{branch systems}, but, as \citet{Sklar1993} and Albert observed, the branches must begin their careers in low-entropy states, and the statistics alone do not supply this. \citet{Penrose1979,Penrose1989} located the specialness of the early universe in its gravitational degrees of freedom, which were nearly uniform, and estimated that the initial state occupies a fraction of about $10^{-10^{123}}$ of the available phase space.

The contemporary debate concerns whether the past hypothesis completes the explanation or calls for one in its turn. \citet{Albert2000} and \citet{Loewer2007,Loewer2012} hold that it forms, with the dynamics and a statistical postulate, the core of a fundamental theory, and Loewer proposes that it is an axiom of the Humean best system. \citet{Callender2004} argues that there is no puzzle: some facts are brute, and a fact about the boundary of the universe is a natural candidate. \citet{Price1996,Price2004} argues that the special initial state is the real puzzle, and that treating it as brute is a double standard. \citet{Earman2006} argues that the hypothesis, as usually formulated, is not even false. We side with Albert, Loewer and Callender, and give a reason, absent from Callender's presentation, why the past hypothesis requires no explanation: it is an axiom of the best system, and the improbability that seems to call for explanation is the product of a misapplied measure.

\section{Framework}\label{s:framework}

\begin{definition}[Macrostate and entropy]\label{d:entropy}
A \emph{macrostate} is a region of the phase space of a system, consisting of the microstates that agree on a chosen set of macroscopic variables, such as density, temperature and pressure in each small cell of the system. The \emph{Boltzmann entropy} of a microstate $x$ is $S_B(x)=k\log|\Gamma_{M(x)}|$, where $M(x)$ is the macrostate containing $x$, $|\Gamma_{M(x)}|$ the phase space volume of that macrostate according to the Liouville measure, and $k$ Boltzmann's constant.
\end{definition}

The Boltzmann entropy is a property of microstates, relative to a partition into macrostates fixed by physically natural macroscopic variables \citep{Goldstein2001}.

\begin{definition}[Temporal asymmetries]\label{d:asym}
The \emph{temporal asymmetries} are the following generalizations, each true in one orientation of the temporal order and false in the other: (a) the \emph{thermodynamic} asymmetry: the entropy of quasi-isolated systems out of equilibrium increases; (b) the \emph{record} asymmetry: present states carry detailed information about states in one orientation and not about states in the other; (c) the \emph{influence} asymmetry: the outcomes of present interventions are correlated, counterfactually, with states in one orientation and not with states in the other; (d) the \emph{radiation} asymmetry: coherent waves diverge from sources and do not converge on sinks.
\end{definition}

\begin{definition}[Direction of time]\label{d:direction}
Time has a \emph{direction} if the temporal asymmetries hold in the same orientation of the temporal order, at almost all places and times. That orientation is the \emph{past-to-future} orientation.
\end{definition}

Whether the earlier-than relation also has a primitive orientation is taken up in Objection~\ref{o:primitive}.

\begin{assumption}[Dynamical symmetry]\label{a:sym}
The fundamental dynamics is invariant under time reversal: if a sequence of microstates is a lawful history, so is the sequence of their time reverses in the reverse order. The time reversal map preserves the Liouville measure and maps each macrostate onto a macrostate of the same entropy.
\end{assumption}

The small violation of time-reversal invariance in the weak interaction is irrelevant to the asymmetries of Definition~\ref{d:asym}, which would be unaffected if it were absent.

\begin{assumption}[Statistical postulate]\label{a:sp}
Given that a system is in macrostate $M$, the probability that its microstate lies in a subset $A$ of $\Gamma_M$ is proportional to the Liouville volume of $A$.
\end{assumption}

\begin{assumption}[Typicality of entropy increase]\label{a:typ}
For the macrostates of systems far from equilibrium, the Liouville measure of the set of microstates in $\Gamma_M$ whose forward evolution does not increase in entropy, until equilibrium is reached, is negligibly small.
\end{assumption}

Assumption~\ref{a:typ} is Boltzmann's conjecture, proved only for idealized systems but granted by every party to the present debate \citep{Goldstein2001}.

\begin{definition}[Past hypothesis]\label{d:ph}
The \emph{past hypothesis} (PH) is the proposition that at one temporal boundary of the universe, the initial boundary, its macrostate $M_0$ is one of very low entropy, of the kind described by the standard cosmological model: matter and radiation close to thermal equilibrium, and spacetime close to homogeneity and isotropy.
\end{definition}

\begin{principle}[Best-system laws]\label{p:bs}
The laws of nature are the axioms and theorems of the best system: the deductive system, stated in natural predicates, that best balances simplicity and informativeness about the total history of the world. Probabilistic axioms are included, and a system's informativeness includes its fit to the actual history.
\end{principle}

Principle~\ref{p:bs} is the Humean account of laws of \citet{Lewis1994}, defended in MET-12-a, where laws are argued to explain without governing.

\begin{principle}[Explanation]\label{p:expl}
To explain a fact is to locate it in a pattern of dependence: to show how it follows, deductively or with high probability, from more general patterns together with particular conditions. A fact \emph{requires further explanation} if and only if it is not already so located by the best available theory, or it is \emph{surprising} relative to that theory: there is an available hypothesis, of comparable simplicity and independent credibility, on which the fact would be much more probable than it is on the theory.
\end{principle}

The first clause is the account of explanation of SCI-03. The second follows \citet{Horwich1982}: a hundred heads in a row calls for explanation because the hypothesis of a two-headed coin is available; a particular mixed series, no less improbable, does not.

\section{The past hypothesis explains the direction of time}\label{s:arrow}

\begin{lemma}[Reversibility]\label{l:rev}
Given Assumptions~\ref{a:sym}--\ref{a:typ}, conditioning only on the present macrostate of a system out of equilibrium yields an overwhelming probability that its entropy was higher in the past, as well as that it will be higher in the future.
\end{lemma}

\begin{proof}
Let $M$ be the present macrostate. By Assumption~\ref{a:typ}, the set $F$ of microstates in $\Gamma_M$ whose forward evolution does not increase in entropy has negligible measure. Let $R$ be the time-reversal map. By Assumption~\ref{a:sym}, $R$ maps $\Gamma_M$ onto the volume $\Gamma_{M^*}$ of a macrostate $M^*$ of equal entropy and preserves measure, and a microstate $x$ had lower entropy in the past if and only if $R(x)$ evolves forward to lower entropy. The set of microstates of $\Gamma_{M}$ that had lower entropy in the past is therefore the image under $R$ of the set of microstates of $\Gamma_{M^*}$ whose forward evolution does not increase in entropy, which has negligible measure by Assumption~\ref{a:typ} applied to $M^*$. By Assumption~\ref{a:sp}, both the forward and the backward exception sets have negligible probability.
\end{proof}

\begin{corollary}[Sceptical catastrophe]\label{c:scep}
Given Assumptions~\ref{a:sym}--\ref{a:typ} and the present macrostate alone, our present records of the past are overwhelmingly likely to be the products of fluctuations rather than of the events they appear to record.
\end{corollary}

\begin{proof}
A record of an event, such as a footprint or a photographic plate, is a present low-entropy configuration whose usual explanation is an interaction with an earlier, still lower-entropy state of the recorded system and the recording device. By Lemma~\ref{l:rev}, the present macrostate alone makes it overwhelmingly probable that entropy was higher in the past, and so that the earlier states which the records seem to record did not occur; the records are then fluctuations.
\end{proof}

Without a further premise, then, statistical mechanics undermines the evidence on which it rests \citep{Albert2000}.

\begin{proposition}[Sufficiency]\label{p:suff}
Given Assumptions~\ref{a:sym}--\ref{a:typ}, conditioning the statistical postulate on PH together with the present macrostate yields an overwhelming probability that entropy has increased from the initial boundary to the present and will continue to increase in quasi-isolated systems out of equilibrium.
\end{proposition}

\begin{proof}
Conditioning on PH restricts the possible microstates of the universe to those whose history passes through $\Gamma_{M_0}$ at the initial boundary. By Assumption~\ref{a:typ}, almost all microstates in $\Gamma_{M_0}$ evolve toward higher entropy; conditioning further on the present macrostate, which those microstates mostly reach, leaves a set of microstates almost all of which have increasing entropy on the way from $M_0$ to the present, since microstates in the present macrostate that came from $M_0$ by an entropy-decreasing route are a negligible fraction of those that came from $M_0$ at all. For the future, Assumption~\ref{a:typ} applies directly to the present macrostate. The same reasoning applies to quasi-isolated subsystems, whose initial states, on becoming isolated, are low-entropy because their environments are: this is Reichenbach's branch-system structure, with the low entropy of the branches inherited from PH.
\end{proof}

\begin{proposition}[Necessity]\label{p:nec}
No proposition consistent with Assumptions~\ref{a:sym}--\ref{a:typ} that constrains only the present and future macrostates of the universe yields the thermodynamic asymmetry toward the past. A boundary condition of low entropy at the past end of the history is required.
\end{proposition}

\begin{proof}
By Lemma~\ref{l:rev}, conditioning on present macrostates yields entropy increase toward the past. A low-entropy future boundary would, by the time reverse of Proposition~\ref{p:suff}, yield entropy decrease toward the future, and a high-entropy one adds nothing. Only a constraint making some past macrostate of lower entropy than the present breaks the symmetry in the right direction, and only one at the initial boundary accounts for the whole observed history.
\end{proof}

The derived asymmetries follow. A record is a system whose present state, together with the assumption that it was in a low-entropy ready state earlier, carries information about an intervening interaction; PH underwrites that assumption toward the past and nothing underwrites it toward the future \citep{Albert2000}. Hence the epistemic asymmetry: the past is known through records, while knowledge of the future is confined, as EPI-25-a argues, to what follows by law from present conditions. The influence asymmetry follows in the manner of \citet{Loewer2007}: given PH, distant macrostates depend counterfactually on small local differences toward the future, and hardly at all toward the past, which is pinned by PH and by the records PH makes reliable. The causal analysis of MET-04, by counterfactual difference-making, inherits this direction, as \citet{Lewis1979} anticipated. The radiation asymmetry follows from the absence of coordinated incoming fields in $M_0$.

\begin{theorem}[The arrow]\label{t:arrow}
Given Assumptions~\ref{a:sym}--\ref{a:typ}, time has a direction because the universe began in a state of low entropy: PH is sufficient, together with the dynamics and the statistical postulate, for the temporal asymmetries of Definition~\ref{d:asym}, and some such boundary condition is necessary for them.
\end{theorem}

\begin{proof}
Sufficiency for the thermodynamic asymmetry is Proposition~\ref{p:suff}, and the record, influence and radiation asymmetries are derived from it in the paragraph preceding the theorem. They hold in the same orientation, that pointing away from the initial boundary, at almost all places and times, so by Definition~\ref{d:direction} time has a direction. Necessity is Proposition~\ref{p:nec}, together with the dependence of the other asymmetries on the thermodynamic one.
\end{proof}

The \enquote{because} of Theorem~\ref{t:arrow} is that of Principle~\ref{p:expl}, not of production: MET-12-a argues that laws do not produce what they summarize.

\section{The past hypothesis requires no further explanation}\label{s:noexpl}

What explains PH? Not something earlier: TSM-14 argues that time had a beginning and that the question of what preceded it is ill-formed, so there is no earlier state to explain the initial one causally. We argue that PH is itself a law, and as such requires no further explanation.

\begin{lemma}[PH is an axiom of the best system]\label{l:axiom}
Given Principle~\ref{p:bs}, PH is, or is a consequence of, an axiom of the best system.
\end{lemma}

\begin{proof}
Compare the system $D$ consisting of the dynamical laws and the statistical postulate with the system $D+\mathrm{PH}$. By Lemma~\ref{l:rev} and Corollary~\ref{c:scep}, $D$ assigns overwhelming probability to entropy increase toward the past and to the unreliability of records; its fit to the actual history, in which entropy has increased from the past and records are reliable, is extremely poor. By Proposition~\ref{p:suff} and Theorem~\ref{t:arrow}, $D+\mathrm{PH}$ assigns high probability to the thermodynamic behavior of every quasi-isolated system in the history of the universe, to the reliability of records, and to the other asymmetries. The gain in informativeness is immense, since these are among the most pervasive regularities in the mosaic. The cost in simplicity is one short statement about the macrostate at one boundary. No system without some such statement achieves comparable fit, by Proposition~\ref{p:nec}. The best system therefore contains PH or a statement entailing it.
\end{proof}

The argument of Lemma~\ref{l:axiom} is that of \citet{Albert2000} and \citet{Loewer2012}, and MET-12-a records that the past hypothesis may be an axiom of the best system. On the Humean account a boundary condition is not disqualified from lawhood by its form: a statement about the initial macrostate earns its place in the best summary of the mosaic by what it summarizes.

\begin{lemma}[Axioms require no further explanation]\label{l:closure}
Given Principles~\ref{p:bs} and~\ref{p:expl}, an axiom of the best system is not left unlocated by the best available theory.
\end{lemma}

\begin{proof}
An axiom of the best system is one of the most general patterns of dependence the best theory contains; if a more general pattern entailed it, a simpler system would state that pattern and derive the axiom. It is located as a summary of the history it is a pattern in, true because the mosaic is as it is, and the mosaic, as MET-12-a argues, is not grounded in anything further. The first clause of the condition for requiring further explanation in Principle~\ref{p:expl} is therefore not met.
\end{proof}

On a governing account of laws, a boundary law would be an anomaly, something a governor had to arrange. On the Humean account, dynamical and boundary laws are two kinds of summary, earning their place by the same criterion and grounded in nothing except the facts they summarize. To ask why the axioms of the best system are true is to ask why the mosaic is as it is, and that question, if it has an answer, is not one that a further physical law could supply. SCI-18 argues that science cannot explain the existence of the conditions that make science possible; PH is among those conditions, since by Corollary~\ref{c:scep} it is what makes records, and therefore evidence about the past, reliable.

It remains to show that PH is not surprising in the sense of Principle~\ref{p:expl}. The case for its being surprising is Penrose's number: PH confines the initial microstate to a region of measure about $10^{-10^{123}}$ of the available phase space, so that if the initial state had been chosen at random, the chance of a state like ours would have been absurdly small \citep{Penrose1989}.

\begin{lemma}[The measure is conditional]\label{l:measure}
The statistical postulate, as it figures in the best system, does not assign a probability to PH.
\end{lemma}

\begin{proof}
In $D+\mathrm{PH}$, the statistical postulate assigns probabilities to microstates conditional on a macrostate (Assumption~\ref{a:sp}), and its role in the system is to supply the distribution over microstates within $\Gamma_{M_0}$ and, derivatively, within the macrostates of later subsystems. The system contains no unconditional probability distribution over the macrostates of the universe at its initial boundary. To compute a probability of PH as the ratio of $|\Gamma_{M_0}|$ to the volume of the whole phase space is to apply the postulate to a condition, namely the whole phase space of the universe, on which the best system does not use it, and whose use would, by Lemma~\ref{l:rev} and Corollary~\ref{c:scep}, worsen the system's fit. That ratio is therefore not a probability that the best system assigns.
\end{proof}

What a measure would make of the initial state if it were applied unconditionally is not evidence about anything unless that measure describes a process by which the initial state was produced, and there was no such process. Our demand for explanation of improbable facts is calibrated on facts within the universe, where the measure is underwritten by the dynamics, and does not transfer to its boundary \citep{Callender2004}. The same diagnosis applies to uniform measures over the constants of nature, which REL-10-a shows cannot be normalized; REL-10-b draws the parallel.

\begin{lemma}[No rival makes PH probable]\label{l:rival}
No available hypothesis of comparable simplicity and independent credibility makes PH much more probable than the best system does.
\end{lemma}

\begin{proof}
We consider the candidates. (i) The fluctuation hypothesis predicts that our records are unreliable (Section~\ref{ss:history}) and so cannot be credible on our evidence. (ii) Inflation does not explain low entropy, since the inflating patch must itself begin in a state of very low entropy; the improbability is relocated, not removed \citep{Penrose1989,AlbrechtSorbo2004}. (iii) Baby-universe proposals such as that of \citet{CarrollChen2004} derive low-entropy beginnings for branches from a high-entropy background, but they depend on unconfirmed features of quantum gravity, and in them the low-entropy beginnings of the branches play the role PH plays here. (iv) Hypotheses on which an agent chose the initial state make it probable only given a specification of the agent's preferences that is not independently credible; REL-10 argues, for the parallel case of fine-tuning, that the claimed improbability is not even defined. None of the candidates satisfies the condition of Principle~\ref{p:expl}.
\end{proof}

\begin{theorem}[No further explanation]\label{t:noexpl}
Given Principles~\ref{p:bs} and~\ref{p:expl} and Assumptions~\ref{a:sym}--\ref{a:typ}, the past hypothesis requires no further explanation.
\end{theorem}

\begin{proof}
By Principle~\ref{p:expl}, PH requires further explanation only if it is not located by the best available theory or it is surprising relative to that theory. By Lemma~\ref{l:axiom}, PH is an axiom of the best system, and by Lemma~\ref{l:closure} it is therefore located. By Lemma~\ref{l:measure}, the best system does not assign it the low probability that Penrose's number expresses, and by Lemma~\ref{l:rival}, no available rival of comparable standing makes it much more probable. It is therefore not surprising.
\end{proof}

\begin{corollary}\label{c:main}
Time has a direction because the universe began in a low-entropy state, and the past hypothesis requires no further explanation.
\end{corollary}

\begin{proof}
Theorems~\ref{t:arrow} and~\ref{t:noexpl}.
\end{proof}

\section{Objections and replies}\label{s:objections}

\begin{objection}[The double standard]\label{o:price}
\citet{Price1996,Price2004} argues that if we found the universe to end in a state of very low entropy, we would regard this as a remarkable fact calling for explanation. To regard the same fact at the beginning as calling for none is to apply a double standard, reading the temporal asymmetry we are trying to explain into our explanatory practice.
\end{objection}

\begin{reply}
We accept the demand for symmetry and meet it. If the universe were found to end in a state of very low entropy, the case would be exactly parallel: a statement of the final macrostate would earn a place in the best system for the same reason that PH does, and by Lemma~\ref{l:closure} it would require no further explanation. Theorem~\ref{t:noexpl} appeals to the place of a boundary condition in the best system and to the absence of a rival on which it would be probable, and both considerations are temporally symmetric. Price is right that there is no temporal asymmetry in what requires explanation; he is wrong to conclude that the low-entropy past requires it.
\end{reply}

\begin{objection}[Not even false]\label{o:earman}
\citet{Earman2006} argues that PH, stated in terms of Boltzmann entropy, is not well defined for the early universe, and that its standard formulations are therefore not even false. An ill-defined proposition can be neither an axiom of the best system nor an explanation of anything.
\end{objection}

\begin{reply}
The objection bears on the formulation of PH, not on its role. Definition~\ref{d:ph} states PH in terms of the macrostate of the standard cosmological model, and the condition that matter was near equilibrium while spacetime was near homogeneity is a well-defined condition on solutions of general relativity; the derivation of the later asymmetries applies the statistical postulate only to later subsystems, where the Boltzmann entropy is well defined. That we do not yet know the exact form of a law does not show that there is no law, any more than uncertainty about the dynamics shows that there are no dynamical laws.
\end{reply}

\begin{objection}[A primitive direction]\label{o:primitive}
\citet{Maudlin2007} holds that time has an intrinsic direction, a primitive feature of the temporal order, and that the thermodynamic asymmetry is a consequence of how matter behaves in a directed time rather than its ground. On this view, time's direction is not explained by PH; PH is stated in terms of a direction it presupposes.
\end{objection}

\begin{reply}
A primitive orientation, if there were one, would not explain why the asymmetries hold: the dynamics are symmetric and a primitive arrow does not enter them, so the asymmetries would still require PH (Proposition~\ref{p:nec}). And the boundary that PH calls initial can be identified without a primitive direction, as the boundary at which entropy is low. A primitive direction is therefore idle in the explanation of the asymmetries. Whether it nonetheless exists is a question about the B-relation that the account need not settle, and TSM-01 and TSM-02 give it no work to do.
\end{reply}

\section{What the result does not establish}\label{s:limits}

The result does not establish the exact content of PH, which awaits a statistical mechanics of gravitating systems. It does not show that PH admits no further explanation, only that it requires none: if a future theory derived PH from simpler axioms, PH would be explained by subsumption, and Lemma~\ref{l:closure} would apply to the new axioms. It does not establish that time flows or that the present is privileged (TSM-02, TSM-15). It does not bear on the openness of the future: ACT-17-a argues that future contingents are unsettled at a moment, and the thermodynamic asymmetry constrains which futures are probable, not which are settled. Finally, it is conditional on the Humean account of laws.

\section{Conclusion}

The dynamical laws are symmetric in time, and the statistical postulate applied to the present macrostate alone predicts entropy increase toward the past as well as the future. A low-entropy boundary condition at the beginning of the universe is necessary and sufficient, together with the dynamics and the statistical postulate, for the thermodynamic asymmetry, and the asymmetries of records, knowledge, influence and radiation derive from it. Time has a direction because the universe began in a low-entropy state. The past hypothesis requires no further explanation: it is an axiom of the best system, whose axioms summarize the mosaic and are not grounded in anything further; the improbability charged against it rests on applying a measure that the best system uses only conditionally on it; and no rival of comparable standing makes it probable.

\begin{thebibliography}{99}\small
\bibitem[Albert(2000)]{Albert2000} Albert, D.~Z. (2000). \emph{Time and Chance}. Harvard University Press.
\bibitem[Albrecht and Sorbo(2004)]{AlbrechtSorbo2004} Albrecht, A., and Sorbo, L. (2004). Can the universe afford inflation? \emph{Physical Review D}, 70, 063528.
\bibitem[Boltzmann(1895)]{Boltzmann1895} Boltzmann, L. (1895). On certain questions of the theory of gases. \emph{Nature}, 51, 413--415.
\bibitem[Boltzmann(1896)]{Boltzmann1896} Boltzmann, L. (1896--1898). \emph{Vorlesungen über Gastheorie}. Barth. English translation by S.~G. Brush, \emph{Lectures on Gas Theory}, University of California Press, 1964.
\bibitem[Callender(2004)]{Callender2004} Callender, C. (2004). There is no puzzle about the low-entropy past. In C.~Hitchcock (ed.), \emph{Contemporary Debates in Philosophy of Science}. Blackwell.
\bibitem[Carroll and Chen(2004)]{CarrollChen2004} Carroll, S.~M., and Chen, J. (2004). Spontaneous inflation and the origin of the arrow of time. arXiv:hep-th/0410270.
\bibitem[Dyson et~al.(2002)]{DysonKlebanSusskind2002} Dyson, L., Kleban, M., and Susskind, L. (2002). Disturbing implications of a cosmological constant. \emph{Journal of High Energy Physics}, 2002(10), 011.
\bibitem[Earman(2006)]{Earman2006} Earman, J. (2006). The \enquote{past hypothesis}: Not even false. \emph{Studies in History and Philosophy of Modern Physics}, 37(3), 399--430.
\bibitem[Eddington(1928)]{Eddington1928} Eddington, A.~S. (1928). \emph{The Nature of the Physical World}. Cambridge University Press.
\bibitem[Feynman(1965)]{Feynman1965} Feynman, R.~P. (1965). \emph{The Character of Physical Law}. MIT Press.
\bibitem[Goldstein(2001)]{Goldstein2001} Goldstein, S. (2001). Boltzmann's approach to statistical mechanics. In J.~Bricmont et~al. (eds.), \emph{Chance in Physics: Foundations and Perspectives}. Springer.
\bibitem[Horwich(1982)]{Horwich1982} Horwich, P. (1982). \emph{Probability and Evidence}. Cambridge University Press.
\bibitem[Lewis(1979)]{Lewis1979} Lewis, D. (1979). Counterfactual dependence and time's arrow. \emph{No\^us}, 13(4), 455--476.
\bibitem[Lewis(1994)]{Lewis1994} Lewis, D. (1994). Humean supervenience debugged. \emph{Mind}, 103(412), 473--490.
\bibitem[Loewer(2007)]{Loewer2007} Loewer, B. (2007). Counterfactuals and the second law. In H.~Price and R.~Corry (eds.), \emph{Causation, Physics, and the Constitution of Reality}. Oxford University Press.
\bibitem[Loewer(2012)]{Loewer2012} Loewer, B. (2012). Two accounts of laws and time. \emph{Philosophical Studies}, 160(1), 115--137.
\bibitem[Loschmidt(1876)]{Loschmidt1876} Loschmidt, J. (1876). Über den Zustand des Wärmegleichgewichtes eines Systems von Körpern mit Rücksicht auf die Schwerkraft. \emph{Sitzungsberichte der Kaiserlichen Akademie der Wissenschaften in Wien, Mathematisch-Naturwissenschaftliche Classe}, 73.
\bibitem[Maudlin(2007)]{Maudlin2007} Maudlin, T. (2007). \emph{The Metaphysics Within Physics}. Oxford University Press.
\bibitem[Penrose(1979)]{Penrose1979} Penrose, R. (1979). Singularities and time-asymmetry. In S.~W. Hawking and W.~Israel (eds.), \emph{General Relativity: An Einstein Centenary Survey}. Cambridge University Press.
\bibitem[Penrose(1989)]{Penrose1989} Penrose, R. (1989). \emph{The Emperor's New Mind}. Oxford University Press.
\bibitem[Price(1996)]{Price1996} Price, H. (1996). \emph{Time's Arrow and Archimedes' Point: New Directions for the Physics of Time}. Oxford University Press.
\bibitem[Price(2004)]{Price2004} Price, H. (2004). On the origins of the arrow of time: Why there is still a puzzle about the low-entropy past. In C.~Hitchcock (ed.), \emph{Contemporary Debates in Philosophy of Science}. Blackwell.
\bibitem[Reichenbach(1956)]{Reichenbach1956} Reichenbach, H. (1956). \emph{The Direction of Time}. University of California Press.
\bibitem[Sklar(1993)]{Sklar1993} Sklar, L. (1993). \emph{Physics and Chance: Philosophical Issues in the Foundations of Statistical Mechanics}. Cambridge University Press.
\end{thebibliography}

\end{document}
